\documentclass[11pt]{article}
\usepackage[margin=1in]{geometry}
\usepackage[T1]{fontenc}
\usepackage{lmodern,amsmath,amssymb,amsthm,mathtools,booktabs,microtype}
\usepackage[colorlinks=true,linkcolor=blue,urlcolor=blue,citecolor=blue]{hyperref}
\newtheorem{theorem}{Theorem}[section]
\newtheorem{proposition}[theorem]{Proposition}
\newtheorem{lemma}[theorem]{Lemma}
\newtheorem{corollary}[theorem]{Corollary}
\theoremstyle{remark}
\newtheorem{remark}[theorem]{Remark}
\newcommand{\Z}{\mathbb Z}
\newcommand{\R}{\mathbb R}
\newcommand{\Q}{\mathbb Q}
\newcommand{\cA}{\mathcal A}
\newcommand{\cM}{\mathcal M}
\newcommand{\la}{\lambda}
\DeclareMathOperator{\cont}{cont}
\DeclareMathOperator{\covol}{covol}
\DeclareMathOperator{\diag}{diag}
\title{Relative arithmetic content in symmetric odd-zeta forms\\
\large Exact reductions, higher-weight decay, and a conditional irrationality criterion}
\author{}
\date{28 October 2026}
\begin{document}
\maketitle
\begin{abstract}
We separate analytic decay from integral normalization in symmetric Pad\'e
forms for odd zeta values. Smith reduction gives an exact formula for the
content of the target and rational rows on the cancellation lattice.
A direct kernel comparison extends the weight-five decay estimate to every
fixed higher weight. An elementary logarithm inequality gives a uniform
exponential margin greater than one third for weights at least seven.
Consequently, nonzero integral forms with subexponential relative arithmetic
cost imply irrationality. The Lean companion proves this conditional
implication, including eventual bounds and arbitrary unbounded indexing,
and an exact integral-projector divisibility criterion.
Signed combinations are also allowed. A target-eliminating cofactor
combination gives an integer lower bound that rules out an overstrong
cofactor-norm saving proposed in the supplied note.
The construction-specific relative-cost bound and the existence of suitable
positive integral cancellation vectors in every higher weight are hypotheses,
not conclusions. This paper does not prove the unconditional irrationality
of every odd zeta value at least seven, or of \(\zeta(5)\).
\end{abstract}

\section{Construction and the arithmetic quantity}
Fix \(s=2r+1\ge5\). Let \(n\ge2r\) be even and select
\(m_j=n-j\), \(0\le j\le r\). Write
\[
 d_n=\operatorname{lcm}(1,\ldots,n),\qquad
 T_{s,n,m}(t)=\frac{(n!)^s}{(m!)^2}
 \frac{\prod_{v=1}^{m}(t-v)(t+n+v)}
      {\prod_{h=0}^{n}(t+h)^s}.
\]
The convergent positive sums are
\[
 \cA_{s,n,m}=\sum_{k=m+1}^{\infty}T_{s,n,m}(k).
\]
Every summand is positive in this range. Partial fractions, reflection
\(t\mapsto-n-t\), and the vanishing total simple-pole residue give
\begin{equation}
 \cA_{s,n,m}=b_{n,m}+
    \sum_{q=1}^{r} a_{q,n,m}\zeta(2q+1).
 \label{eq:expansion}
\end{equation}
The coefficients are rational. Weighted Taylor expansion at each pole
shows
\[
 d_n^{s-(2q+1)}a_{q,n,m}\in\Z,\qquad d_n^s b_{n,m}\in\Z.
\]
For completeness, factor \(T=P_mQ_n^s\), where
\[
 P_m(t)=\binom{t-1}{m}\binom{t+n+m}{m},\qquad
 Q_n(t)=\frac{n!}{\prod_{h=0}^{n}(t+h)}.
\]
Vandermonde expansion shows that \(d_n^\ell\) clears the coefficient
of \(x^\ell\) in \(P_m(-j+x)\). The same statement holds for
\([xQ_n(-j+x)]^s\), since its factors are
\((1+x/(h-j))^{-s}\), apart from an integral leading coefficient.
Thus \(d_n^{s-k}\) clears each coefficient of \((t+j)^{-k}\).
Summing the partial fractions proves the displayed denominator bounds.

Clear all rows by \(D_n=d_n^s\). Let \(C_{s,n}\) be the
\((r-1)\)-by-\((r+1)\) integer matrix of coefficients of the lower
periods \(\zeta(3),\ldots,\zeta(s-2)\), and let \(Y_{s,n},Z_{s,n}\)
be the cleared target and constant rows. Whenever \(C_{s,n}\) has full
row rank, its integral kernel
\[
 \Lambda_{s,n}=\ker_{\Z} C_{s,n}
\]
is a saturated lattice of rank two. For an integral basis \(a,b\), define
\begin{equation}
 \Gamma_{s,n}=\gcd\bigl(Y a,Y b,Z a,Z b\bigr).
 \label{eq:gamma}
\end{equation}
Here and below gcds are nonnegative; we require \(\Gamma_{s,n}>0\)
when dividing by it. This gcd is independent of the lattice basis.

\begin{proposition}[Kernel rank and saturation]
If \(\operatorname{rank}_{\Q}C=r-1\), then
\(\ker_{\Z}C\cong\Z^2\), and this kernel is saturated in \(\Z^{r+1}\).
\end{proposition}
\begin{proof}
Smith normal form supplies unimodular matrices \(U,V\) with
\[
 UCV=(\diag(d_1,\ldots,d_{r-1})\mid0\mid0),\qquad d_i>0.
\]
The equation \(Cv=0\) says precisely that the first \(r-1\) coordinates
of \(V^{-1}v\) are zero. The final two coordinates are arbitrary integers,
so the final two columns of \(V\) give an integral kernel basis.
If \(kv\) belongs to this kernel, with \(k\in\Z\setminus\{0\}\), then
\(0=C(kv)=kCv\). Torsion-freeness of \(\Z^{r-1}\) gives \(Cv=0\),
which proves saturation.
\end{proof}

At weight five the established adjacent-sign theorem gives
\(x_{0,n}x_{1,n}<0\) for the single lower-period row at every even
\(n\ge2\). That row is nonzero and therefore has rational rank one;
its integral kernel has rank two by the proposition.
For higher weights, counting the \(r-1\) rows is not a proof of their
independence. The supplied rank argument explicitly retains that hypothesis.

\begin{proposition}[Exact content by Smith reduction]
Let \(\Delta_k(B)\) be the nonnegative gcd of all \(k\)-minors of an
integer matrix \(B\), with \(\Delta_0(B)=1\).
For a full-row-rank matrix \(C\) with \(r-1\) rows and \(r+1\) columns,
\begin{equation}
 \Gamma=
 \frac{\gcd\bigl(\Delta_r(C;Y),\Delta_r(C;Z)\bigr)}
      {\Delta_{r-1}(C)}.
 \label{eq:smith}
\end{equation}
The notation \((C;Y)\) means append the row \(Y\).
\end{proposition}
\begin{proof}
Choose unimodular \(U,V\) with
\(UCV=(\diag(d_1,\ldots,d_{r-1})\mid0)\).
The last two columns of \(V\) are a basis of the integral kernel.
The gcd of the maximal minors of \(C\) is \(\prod_i d_i\).
The only possibly nonzero \(r\)-minors after appending \(Y\) are
\((\prod_i d_i)Ya\) and \((\prod_i d_i)Yb\), up to sign.
Unimodular row and column transformations preserve determinantal
divisors. Repeat with \(Z\) and take the gcd.
\end{proof}
Consequently, with the convention \(v_p(0)=+\infty\),
\begin{equation}
 v_p(\Gamma)=
 \min\{v_p(\Delta_r(C;Y)),v_p(\Delta_r(C;Z))\}
       -v_p(\Delta_{r-1}(C)).
 \label{eq:valuation}
\end{equation}
This is a divisibility problem, distinct from the real growth of the minors.

\section{The exact integer-form bound}
Suppose \(w_{s,n}\in\Lambda_{s,n}\cap\Z_{\ge0}^{r+1}\) is nonzero.
Put \(M_{s,n}=\|w_{s,n}\|_1\) and
\[
 h(w)=\gcd(Yw,Zw),\qquad
 L_{s,n}=\frac{(Yw)\zeta(s)+Zw}{h(w)}.
\]
The numerator equals \(D_n\sum_j w_j\cA_{s,n,n-j}>0\).
In particular \(h(w)>0\). Since \(\Gamma_{s,n}\mid h(w)\), we obtain
\begin{equation}
 0<L_{s,n}\le
 d_n^s\frac{M_{s,n}}{\Gamma_{s,n}}\,
 S_{s,n},\qquad
 S_{s,n}=\max_{0\le j\le r}\cA_{s,n,n-j}.
 \label{eq:bound}
\end{equation}
Both coefficients of \(L_{s,n}\) are integers.
Rational normalization of a kernel vector to mass one does not prove
a bound for \(M_{s,n}/\Gamma_{s,n}\), since clearing its denominator
must still be included in \eqref{eq:bound}.

\section{A direct analytic extension to higher weights}
The following argument makes the weight dependence explicit without a
new saddle computation at each weight.

\begin{lemma}[Extra powers of the kernel]
For \(n\ge2r\), \(0\le j\le r\), and \(k\ge n-j+1\),
\[
 Q_n(k)\le3^j4^{-n},\qquad
 T_{s,n,n-j}(k)=Q_n(k)^{s-5}T_{5,n,n-j}(k).
\]
\end{lemma}
\begin{proof}
The binomial sum and maximality of the central binomial coefficient give
\(4^n\le(2n+1)\binom{2n}{n}\), hence
\[
 Q_n(n+1)=\frac{1}{(2n+1)\binom{2n}{n}}\le4^{-n}.
\]
The function \(Q_n(k)\) decreases for positive \(k\).
Also \(Q_n(k)/Q_n(k+1)=(k+n+1)/k\).
For \(k\ge n-j+1\), \(n\ge2j\), this ratio is at most three.
At most \(j\) backward steps from \(n+1\) suffice.
The identity follows by cancellation of the definitions.
\end{proof}

For comparison of the fixed shifts at weight five, exact factorial
cancellation gives
\begin{align}
 \frac{T_{5,n,n-j}(k)}{T_{5,n,n}(k+j)}
 &=
 \left(\frac{n!}{(n-j)!}\right)^2
 \prod_{h=1}^j\left(\frac{k+n+h}{k+h-1}\right)^6
 \left(\prod_{h=1}^{2j}(k+2n-j+h)\right)^{-1} \notag\\
 &\le (3^6/4)^j. \label{eq:shift}
\end{align}
Indeed the first factorial ratio is at most \(n^j\), each middle ratio
is at most three, and each final factor is at least \(2n\).
The change \(k\mapsto k+j\) carries the summation range onto \(k+j\ge n+1\).

The elementary weight-five estimate, proved in Appendix A, is
\begin{equation}
 \cA_{5,n,n}\le C_0(n+1)^{12}140^{-n},
 \qquad C_0=12\cdot3^{12}\cdot5^{10}e^{10}.
 \label{eq:base}
\end{equation}
It is not a Lean theorem in the present companion. Appendix A also proves
the sharp bound \(\cA_{5,n,n}\le C_0(n+1)^{12}e^{\la_5n}\).
The derivation locates
the unique maximum of
\[
 6[x\log x-(x+1)\log(x+1)]
 +(x+2)\log(x+2)-(x-1)\log(x-1).
\]
Writing \(y=x(x+1)\), its stationary equation is
\(3y^3-4y^2-5y-1=0\).
The maximum occurs below \(1057/1000\), where the resulting exponential
base is below \(1/140\). The complete summation and tail estimates in
Appendix A supply the prefactor in \eqref{eq:base} and certify the sharper
interval \(-4.970794<\la_5<-4.970793\).

\begin{proposition}[Higher-weight raw decay]
For each fixed odd \(s=2r+1\ge7\) and
even \(n\ge2r\),
\begin{equation}
 S_{s,n}\le C_s(n+1)^{12}\exp(-a_s n),\qquad
 a_s=\log140+(s-5)\log4,\qquad
 C_s=C_0\,3^{r(s+1)}.
 \label{eq:raw}
\end{equation}
\end{proposition}
\begin{proof}
Apply the lemma termwise, sum, and use \eqref{eq:shift}.
For the shift \(j\), the constant obtained is
\(3^{j(s-5)}(3^6/4)^j\le3^{r(s+1)}\).
The powers \(140^{-n}4^{-(s-5)n}\) equal \(\exp(-a_s n)\).
\end{proof}

\begin{lemma}[Explicit margin, checked in Lean]
For every integer \(s\ge7\),
\begin{equation}
 a_s>s+\tfrac13.
 \label{eq:margin}
\end{equation}
\end{lemma}
\begin{proof}
The elementary inequality \(2x/(x+2)<\log(1+x)\), applied at \(x=1\),
gives \(\log2>2/3\). Also
\(\log140>7\log2>14/3\) and \(\log4=2\log2>4/3\).
Thus
\[
 a_s>\frac{14}{3}+\frac43(s-5)=\frac{4s-6}{3}
 \ge s+\frac13.
\]
\end{proof}

\section{The precise conditional irrationality theorem}
All limits in this section may be restricted to an unbounded sequence
of eligible even indices. Write \(\log^+x=\max(0,\log x)\) for \(x>0\).

\begin{theorem}[Relative arithmetic cost]
Fix an odd \(s\ge7\). Suppose that the full-rank cancellation setup,
the positive integral vectors, and the integer forms in
\eqref{eq:bound} exist along an unbounded sequence of even indices.
Suppose \eqref{eq:raw} holds and
\begin{equation}
 \log^+(M_{s,n}/\Gamma_{s,n})=o(n).
 \label{eq:cost}
\end{equation}
Then \(\zeta(s)\) is irrational.
\end{theorem}
\begin{proof}
The prime number theorem yields \(\log d_n=n+o(n)\).
For fixed \(s\), the denominator term has upper exponential rate \(s\).
The polynomial factor in \eqref{eq:raw} has rate zero.
By \eqref{eq:cost} the relative arithmetic cost also has upper rate
at most zero. Taking logarithms in \eqref{eq:bound} therefore gives
\[
 \limsup\frac{\log L_{s,n}}n\le s-a_s<-\tfrac13.
\]
Hence the nonzero integer linear forms tend to zero.
If \(\zeta(s)=p/q\), \(p,q\in\Z\), \(q\ne0\), then
\(qL_{s,n}\) is a nonzero integer and has absolute value at least one,
contradicting convergence to zero.
\end{proof}

For clarity, the Lean theorem uses eventual exponential envelopes
instead of a quotient of logarithms. With \(D_n=d_n^s\),
\(W_n=M_{s,n}/\Gamma_{s,n}\), its assumptions are, for every
\(\varepsilon>0\), eventually
\[
 D_n\le e^{(s+\varepsilon)n},\qquad
 W_n\le e^{\varepsilon n},\qquad
 S_n\le e^{(-a_s+\varepsilon)n}.
\]
These are supplied as hypotheses to the theorem, together with
\eqref{eq:bound} and arbitrarily late nonvanishing. Choosing
\(\varepsilon=1/12\) proves the explicit eventual bound
\[
 |L_n|\le e^{-n/12}\longrightarrow0.
\]
The formal theorem permits any index map tending to infinity, including
the even subsequence; no finite-index extrapolation is used.

\section{An integral projector for the divisibility obligation}
At weight five the cancellation row is a nonzero integer row \(x\).
Put \(g=\gcd(x_1,x_2,x_3)\), \(\ell=x/g\), so that \(\ell\) is primitive.
Choose \(u\in\Z^3\) with \(\ell u=1\). Define
\[
 \Pi=I-u\ell.
\]
Then \(\Pi^2=\Pi\), \(\ell\Pi=0\), and
\(\Pi v=v\) whenever \(\ell v=0\). Its integral image is exactly
\(\Lambda=\ker_\Z\ell\).

\begin{proposition}[Universal congruence test, checked in Lean]
For any integer matrix \(F\) with three columns and any \(c\in\Z\),
\begin{equation}
 c\mid(F\Pi)_{ij}\ \text{for every }i,j
 \quad\Longleftrightarrow\quad
 c\mid(Fv)_i\ \text{for every }v\in\Lambda\text{ and every }i.
 \label{eq:projector}
\end{equation}
The gcd of the entries of \(F\Pi\) is independent of the choice of \(u\),
up to multiplication by a unit.
\end{proposition}
\begin{proof}
If the left side holds, write \(v=\Pi v\) and expand \(F\Pi v\)
as an integral linear combination of its columns.
Conversely apply the right side to each column \(\Pi e_j\), which
lies in \(\Lambda\). The equivalence characterizes all divisors of
the entry gcd and is independent of the section \(u\).
\end{proof}
Taking \(F\) to have rows \(Y,Z\) identifies its entry gcd with
\(\Gamma\) in \eqref{eq:gamma}. This formulation accounts for division
by \(g\) before imposing higher-order congruences. For \(c=p^e\), it
is an exact reformulation of the local obligation, not a lower bound
for its exponent.

\section{Shape control and the weight-five boundary}
The augmented-rank condition used here has the following exact proof.
Choose an integral kernel basis \(a,b\), under the lower-row rank condition,
and put
\[
 B=\begin{pmatrix}Y(a)&Y(b)\\ Z(a)&Z(b)\end{pmatrix}.
\]
Apply the Smith transformations from Section 1 to the augmented matrix.
Its first rows have diagonal entries \(d_1,\ldots,d_{r-1}\), and its
lower right corner is \(B\). Unimodular transformations change determinant
only by a sign. Expansion along the first rows proves
\begin{equation}
 |\det(C;Y;Z)|=\Delta_{r-1}(C)|\det B|.
 \label{eq:augmented}
\end{equation}
Thus the augmented square matrix has full rational rank if and only if
\(\det B\ne0\). This proves the equivalence; it does not assert determinant
nonvanishing for every member of the higher-weight Pad\'e family.

Let \(a,b\) be a Gauss-reduced basis of a saturated rank-two cancellation
lattice, \(\la_1=\|a\|\le\la_2=\|b\|\), and
\(H=\|a\wedge b\|\). Set
\[
 U_Y=Y(a)b-Y(b)a,\qquad
 U_Z=Z(a)b-Z(b)a,\qquad
 \cM=\max(\|U_Y\|,\|U_Z\|).
\]
Assume the augmented square matrix \((C;Y;Z)\) has nonzero determinant.

\begin{proposition}[Cofactor comparison]
Under these assumptions,
\[
 \la_2\le\frac{\cM}{\Gamma},\qquad
 \kappa=\frac{\la_2}{\la_1}
       \le\frac{\cM^2}{\Gamma^2 H},\qquad
 \frac{\Gamma}{\la_2}\ge\frac{\Gamma^2}{\cM}.
\]
\end{proposition}
\begin{proof}
The vectors \(U_Y/\Gamma,U_Z/\Gamma\) are independent lattice vectors.
The second successive minimum is therefore at most
\(\cM/\Gamma\); for a Gauss-reduced rank-two basis it equals \(\la_2\).
Since \(H=\la_1\la_2\sin\theta\) and \(\sin\theta\le1\),
\(\kappa\le\la_2^2/H\). Combine the two inequalities.
\end{proof}

In weight five, positive rounding gives the separate sufficient bound
\[
 M_{5,n}\le3(1+2R_n)\la_{2,n}.
\]
The row-balance result \(\log R_n=o(n)\) removes this factor from
the exponential budget, but not the factor \(\la_{2,n}/\Gamma_{5,n}\).
Using only the explicit base \(140\), a sufficient joint estimate is
\[
 \Gamma_{5,n}\ge e^{0.06n}\la_{2,n}.
\]
The sharper saddle value
\(-4.970794<\la_5<-4.970793\), certified in Appendix A,
reduces \(0.06\) to \(0.03\).
The sharper constant enclosure and this sufficient joint estimate are
not formalized in the new Lean modules. Neither joint arithmetic
estimate has been proved for the constructed family.
In particular this section supplies no unconditional theorem for
\(\zeta(5)\).

\section{Signed cancellation and a cofactor restriction}
For real weights of either sign, the triangle inequality gives
\begin{equation}
 \left|\sum_j w_j\cA_j\right|
       \le\left(\sum_j|w_j|\right)\max_j|\cA_j|.
 \label{eq:signed}
\end{equation}
Thus nonnegativity is not essential to the upper estimate. Its role in
Section 2 is to guarantee nonvanishing. A different guarantee is available
for two coefficient pairs: if \(ad-bc\ne0\), the forms
\(a\alpha+c\) and \(b\alpha+d\) cannot both vanish, since
\[
 a(b\alpha+d)-b(a\alpha+c)=ad-bc.
\]
One may therefore work with a pair of independent signed forms, provided
both satisfy the required arithmetic bounds. These two elementary
assertions are proved in the signed Lean companion.

However, a signed combination that eliminates the target itself is a
nonzero rational period and has an integral lower bound after clearing.
In the notation of Section 6, put
\[
 \delta=Y(a)Z(b)-Y(b)Z(a)\ne0.
\]
Then \(C U_Y=0\), \(Y(U_Y)=0\), and \(Z(U_Y)=\delta\).
Since \(\Gamma\) divides all four displayed evaluations,
\(\Gamma^2\mid\delta\), so \(|\delta|\ge\Gamma^2\).
Applying the cleared period identity and \eqref{eq:signed} gives
\begin{equation}
 \boxed{\Gamma_{s,n}^2
  \le |\delta|
  =d_n^s\left|\sum_j(U_Y)_j\cA_{s,n,n-j}\right|
  \le d_n^s\sqrt{r+1}\,\cM_{s,n}S_{s,n}.}
 \label{eq:cofactorlower}
\end{equation}
Here \(\|U_Y\|_1\le\sqrt{r+1}\|U_Y\|_2\le\sqrt{r+1}\cM\).
The determinant divisibility, signed bound, and integer-divisor lower
bound are proved in \texttt{SignedCancellation.lean}; their application
to these coefficient rows is the displayed paper derivation.

It follows from \eqref{eq:raw} and \(\log d_n=n+o(n)\) that
\[
 \limsup\frac1n\log\frac{\Gamma_{s,n}^2}{\cM_{s,n}}
       \le s-a_s.
\]
In particular, for \(s\ge7\) this upper rate is less than \(-1/3\).
At weight five the coarse base \(140\) gives the upper rate
\(5-\log140<0.058358\); the sharper saddle rate gives
\(5+\la_5<0.029207\).
Consequently the sufficient cofactor strengthenings
\(\Gamma_{5,n}^2\ge e^{0.06n}\cM_{5,n}\), or
\(\Gamma_{5,n}^2\ge e^{0.03n}\cM_{5,n}\) with the sharp saddle bound,
cannot hold along an unbounded sequence with nonzero augmented determinant.
This restriction does not invalidate the comparison in Section 6.
It shows that using its cofactor upper bound in place of the actual
second minimum is too strong for that proposed arithmetic target.
The weaker bound for a selected primitive form remains a separate problem.

\section{Row balance and positive integral cost}
The relative-cost estimate needs arithmetic information about the specified
family. In particular it does not follow from row balance and the existence
of a positive integral kernel vector alone. This can be proved exactly,
without extrapolating from numerical examples.

For an integer \(t\ge2\), take
\[
 C=(t,-(t-1),t),\qquad Y=(1,0,0),\qquad Z=(0,0,1).
\]
The cancellation row is primitive, has full row rank, and its positive to
negative row-sum ratio is \(2t/(t-1)\le4\).
The integral kernel contains \((1,0,-1)\), on which \(Y\) evaluates to one,
so its target-and-constant content is \(\Gamma=1\).
If \(w=(a,b,c)\ge0\) is a nonzero integral kernel vector, then
\[
 t(a+c)=(t-1)b,\qquad k=b-a-c\in\Z_{\ge1}.
\]
It follows that \(a+c=(t-1)k\), \(b=tk\), and
\[
 \|w\|_1=(2t-1)k\ge2t-1.
\]
Equality is attained by \((t-1,t,0)\).
Choosing \(t=2^n\) gives a bounded row-balance ratio but forces every
nonzero nonnegative cancellation vector to have exponentially growing
\(M/\Gamma\). The lower bound, its attaining vector, the row-ratio bound,
and the unit-content implication are checked in Lean.
This family is not the odd-zeta family: it proves a non-implication between
the general structural properties and the desired arithmetic estimate.

\section{Formalization and reproducibility}
The archive contains the pinned Lean project and every source module.
The three additions are:
\begin{itemize}
\item \texttt{OddZeta/HighWeightRelativeCost.lean}:
the logarithm margin, convergence from eventual rate bounds, and the
conditional irrationality theorem for \texttt{oddZetaValue k},
\(k\ge3\), meaning the real part of \(\zeta(2k+1)\).
\item \texttt{OddZeta/KernelProjectorContent.lean}:
the integral projection, idempotence, universal divisibility equivalence,
and independence of the content from the section.
\item \texttt{OddZeta/PositiveCancellationCost.lean}:
the exact minimal positive mass calculation above, with the row-ratio and
unit-content assertions.
\end{itemize}
The publication folder also contains \texttt{SignedCancellation.lean}
and exact cancellation examples at weight seven and indices six and eight.
Two independent rational coefficient algorithms agree at these indices,
and Lean verifies the four resulting integer cancellation identities.
They confirm actual cancellation, with no extrapolation to all indices.
All parameters of the irrationality criterion remain explicit.
In particular the cost bound and nonvanishing are hypotheses, and no
axiom asserts their existence for the Pad\'e family.
The complete source declaration, rather than its theorem name alone,
specifies the mathematical assertion checked by Lean.

A successful Lean compilation verifies these formal statements relative
to the pinned library. It does not verify paper-only derivations such
as \eqref{eq:smith}, \eqref{eq:base}, or the Gauss-reduction comparison.
The accompanying verification report distinguishes these scopes.
Rebuild with the pinned toolchain, using
\texttt{lake exe cache get} and \texttt{lake build}.
The provided lock file records the library revisions.

An unconditional application to all odd weights would require a proof
of \eqref{eq:cost} for explicitly constructed nonnegative integral
cancellation vectors at each fixed weight. The exact projector criterion
and the Smith formula specify the arithmetic content entering that task.

\appendix
\section{Global decay with the sharp saddle constant}
Here is the full analytic proof of \eqref{eq:base} and its sharper version.
For \(x>1\), define
\[
 \Phi(x)=6[x\log x-(x+1)\log(x+1)]
       +(x+2)\log(x+2)-(x-1)\log(x-1).
\]
Its derivative has the sign of
\(x^6(x+2)-(x-1)(x+1)^6\). With \(y=x(x+1)\), the negative of this
expression is \(p(y)=3y^3-4y^2-5y-1\).
For \(y\ge2\), \(p'(y)=9y^2-8y-5>0\), while \(p(2)=-3\).
Therefore \(\Phi\) increases to a unique \(x_*>1\) and decreases thereafter.
Subtracting \(x_*\Phi'(x_*)=0\) from its value gives
\[
 \la_5=\Phi(x_*)=\log q(x_*),\qquad
 q(x)=\frac{(x+2)^2(x-1)}{(x+1)^6}.
\]
Set \(a=1057/1000\). Exact integer arithmetic proves
\begin{align*}
 1000^6p(a(a+1))&=54685013669880747>0,\\
 2057^6-140\cdot3057^2\cdot57\cdot1000^3
       &=1178944474661131249>0.
\end{align*}
Hence \(1<x_*<a\). On this interval
\[
 (\log q)'=\frac2{x+2}+\frac1{x-1}-\frac6{x+1}
           >\frac{1000}{57}-3>0.
\]
The second integer inequality is exactly \(q(a)<1/140\), so
\(\la_5<-\log140\) is a global bound.

For the sharp certificate put
\[
 x_- =\frac{1056126503577}{10^{12}},\qquad
 x_+ =\frac{1056126503578}{10^{12}}.
\]
Exact rational evaluation gives \(p(x_-(x_-+1))<0\) and
\(p(x_+(x_++1))>0\). Thus \(x_-<x_*<x_+\), and monotonicity
of \(q\) encloses \(\la_5\) between the logarithms of these rational values.
For \(1\le u\le2\), let \(z=(u-1)/(u+1)\). The convergent series gives
\[
 0\le\log u-2\sum_{k=0}^{L-1}\frac{z^{2k+1}}{2k+1}
 \le\frac{2z^{2L+1}}{(2L+1)(1-z^2)}.
\]
For any positive rational argument, first factor out an integral power of
two. Applying this formula with \(L=32\) and outward rational rounding
proves
\[
 -4.970793860358<\la_5<-4.970793860341.
\]
The signs and rational enclosures are reproduced by
\texttt{check\_high\_weight\_rates.py}. They certify a single global maximum,
without extrapolation from indices.

To bound the infinite sum, write \(T=T_{5,n,n}\). For \(t\ge2n\),
logarithmic differentiation gives
\[
 \frac{d}{dt}\log(t^2T(t))\le-\frac{E}{t}+\frac{B}{t^2},
 \qquad E=3n+3,\quad B=7n^2+5n.
\]
Indeed apply \(1/(t-v)\le1/t+2v/t^2\),
\(1/(t+n+v)\le1/t\), and \(1/(t+h)\ge1/t-h/t^2\)
to its finite logarithmic derivative. Set
\(K=\max(2n,n+1,\lceil B/E\rceil)\le3n\).
Then \(t^2T(t)\) is decreasing for \(t\ge K\), so
\(T(k)\le T(K)(K/k)^2\) there. The maximum summand \(P_n\)
is attained at some \(n+1\le k\le K\). An integral comparison of the
tail of \(\sum k^{-2}\) gives
\[
 \cA_{5,n,n}\le(2K+1)P_n\le(6n+1)P_n.
\]

For \(f(v)=v\log v-v\), \(f(0)=0\), integral comparison also gives
\[
 f(h)\le\log(h!)\le f(h)+1+\log(h+1)\qquad(h\in\Z_{\ge0}).
\]
At a maximizing index \(n+1\le k\le3n\), the exact factorial identity is
\[
 T(k)=\frac{(n!)^3((k-1)!)^6(k+2n)!}
             {(k-n-1)!((k+n)!)^6}.
\]
Apply the upper factorial bound to the ten numerator factors and the lower
bound to the denominator. The total error is at most
\(10[1+\log(5n+1)]\). In the main terms,
\(f(k-1)-f(k)\le0\) and
\(f(k-n)-f(k-n-1)\le\log(k-n)\le\log(2n)\), including \(k-n=1\).
Replacing these arguments accordingly leaves exactly
\(n\Phi(k/n)\le n\la_5\), apart from error at most \(\log(2n)\).
Consequently, for every even \(n\ge4\),
\begin{align*}
 \cA_{5,n,n}
 &\le2e^{10}n(6n+1)(5n+1)^{10}e^{\la_5n}\\
 &\le12\cdot5^{10}e^{10}(n+1)^{12}e^{\la_5n}
 \le C_0(n+1)^{12}140^{-n}.
\end{align*}
This proves the complete raw decay estimate. Section 3 extends it to every
higher fixed weight. The primitive integer normalization cost remains in
\eqref{eq:bound}; no analytic step here discards that factor.

\begin{thebibliography}{9}
\bibitem{smithnote}
\emph{A Smith-content reduction for odd-zeta Pad\'e forms},
user-supplied technical note, 28 October 2026, five pages.

\bibitem{collected}
\emph{Odd-zeta collected notes},
user-supplied compilation of six working notes, thirty pages.

\bibitem{repository}
Companion Lean project and derivation notes,
\href{https://github.com/ranbirsurii/irrationality-perhaps-}{GitHub repository}.
The archive contains a source snapshot and a verification report.
\end{thebibliography}
\end{document}
